\section{Simulation Results}
\label{sec:sim-results-full}

This section gives the full results summarized in Section~\ref{sec:sim-results} of the paper; Fig.~\ref{fig:sim-sens} is in the paper.

\noindent\fbox{\parbox{0.96\columnwidth}{\small All values in this section are \textsc{synthetic} unless marked \textsc{actual} or stated as \textsc{assumed} inputs; rows marked $^{\ddagger}$ are \textsc{proposed} components simulated with assumed rates. The synthetic values come from the corrected seeded simulation (v2, post hoc) of Section~\ref{sec:sim}, not from running the agent. They show the consequences of the stated assumptions and the measured component behaviour; they do not establish how the deployed agent performs. Every $p$-value and interval here is conditional on those assumptions.}}

\subsection{Baselines and Verification (Studies A and B; H1, H2)}

\begin{table*}[t]
\centering
\caption{\textsc{Simulation-based (synthetic)} Study~B: completion decisions of four configurations on the same 100 held-out synthetic scenarios. Confusion counts and rates with Wilson 95\% intervals. $^{\ddagger}$Not implemented (proposed).}
\label{tab:sim-baseline}
\footnotesize
\setlength{\tabcolsep}{5pt}
\begin{tabular}{@{}lccccccccc@{}}
\toprule
Configuration & TP & FP & FN & TN & TSR & FCR & FRR & F1 & Acc. \\
\midrule
No verification & 43 & 28 & 3 & 26 & 43/100 [0.34,0.53] & 28/71 [0.29,0.51] & 3/46 [0.02,0.18] & .735 & .690 \\
Basic success detection & 42 & 24 & 4 & 30 & 42/100 [0.33,0.52] & 24/66 [0.26,0.48] & 4/46 [0.03,0.20] & .750 & .720 \\
Evidence-gated ($\Phi$) & 37 & 19 & 10 & 34 & 37/100 [0.28,0.47] & 19/56 [0.23,0.47] & 10/47 [0.12,0.35] & .718 & .710 \\
Enhanced validation$^{\ddagger}$ & 39 & 8 & 12 & 41 & 39/100 [0.30,0.49] & 8/47 [0.09,0.30] & 12/51 [0.14,0.37] & .796 & .800 \\
\bottomrule
\end{tabular}
\end{table*}


\begin{table}[t]
\centering
\caption{\textsc{Simulation-based (synthetic)} Paired comparisons with the evidence-gated configuration, exact McNemar tests on discordant pairs $b/c$ (first count: only the other configuration; second: only $\Phi$). FC: episode ends in a false completion (100 scenarios). Decision: first decision correct (71 claims). Success: task achieved and reported (100 scenarios). MC: mean FCR and 2.5--97.5th percentiles over 1000 regenerated scenario sets. Synthetic $p$-values and intervals are conditional on the simulation assumptions; they are not empirical tests of the agent.}
\label{tab:sim-tests}
\scriptsize
\setlength{\tabcolsep}{1.8pt}
\begin{tabular}{@{}lcccccc@{}}
\toprule
Configuration & FC $b/c$ & $p$ & Decision $b/c$ & $p$ & Success $b/c$; $p$ & MC FCR \\
\midrule
No verification & 9/0 & 0.004 & 7/9 & 0.804 & 7/1; 0.070 & 0.36 [0.25,0.48] \\
Basic success detection & 5/0 & 0.062 & 6/5 & 1.000 & 6/1; 0.125 & 0.33 [0.22,0.44] \\
Evidence-gated ($\Phi$) & ref. & ref. & ref. & ref. & ref. & 0.25 [0.15,0.37] \\
Enhanced validation$^{\ddagger}$ & 0/11 & 0.001 & 11/2 & 0.022 & 4/2; 0.688 & 0.14 [0.05,0.24] \\
\bottomrule
\end{tabular}
\end{table}


Table~\ref{tab:sim-baseline} compares four completion rules on the same 100 scenarios of the reported synthetic set (seed 202) at two levels that must not be mixed. At the \emph{decision level}, the planner claims completion in 74 scenarios, 46 of them truly complete at that moment; the four verifiers face the same 74 claims and the same truth, so their decisions are paired. At the \emph{policy level}, an episode's final state may differ between configurations for the same scenario, because a rejected claim returns control to planning and the episode may then finish (or, after a false rejection, never be reported); task success and episode-level false completions are paired outcomes of different policies, not of one ground truth.

\emph{H1 (false completions).} Episodes end in a false completion 19 times with $\Phi$ against 28 without verification (FCR 0.30 against 0.38; discordant pairs 9 and 0; Table~\ref{tab:sim-tests}); basic success detection ends in 24. In the model the direction is fixed by construction (in the code, a continuation after a false rejection could later be accepted on a changed page, which the model omits): $\Phi$ can only withhold a claim that would otherwise be accepted, so no scenario can end in a false completion with $\Phi$ alone, one discordant cell is always 0, and only the size of the reduction is informative. Across 1{,}000 regenerated scenario sets, with the derived rates drawn from their posteriors, the reduction averages 0.11 false completions per scenario (2.5--97.5th percentiles 0.05--0.18); because one discordant cell is structurally 0, a McNemar $p$ would only restate the 9 discordant scenarios, so we report this interval instead. Under the stated assumptions the simulation is consistent with H1; it is not empirical confirmation.

\emph{H2 (decision reliability).} On the 74 claims, $\Phi$ decides correctly in 52 (accuracy 0.70, Wilson 0.59--0.79) and accepting every claim in 46 (0.62, 0.51--0.72): $\Phi$ is correct alone in 9 claims, all rejections of premature claims, and wrong alone in 3, all false rejections (two of them on scenarios with missing evidence, which $\Phi$ rejects at an \textsc{assumed} rate of 0.8; the measured false-rejection rate is 0 of 24) (exact McNemar $p=0.15$). The difference has the hypothesized direction but is not significant in this set, and across regenerated sets $\Phi$ is better at $p<0.05$ in only 42\% (mean accuracy 0.74 against 0.64), so H2 is not supported for $\Phi$ at this sample size (74 claims). The proposed enhanced validation reaches 61 of 74 (0.82; against $\Phi$ 11 and 2, $p=0.022$), a difference that follows from its assumed detection rates. At the policy level, the simulated $\Phi$ configuration ends with 44 of 100 synthetic scenarios achieved and reported, against 46 without verification (3 only without verification, 1 only with $\Phi$; $p=0.63$): in the simulation, false rejections cost 3 successes and continuation after a correct rejection recovers 1.

\subsection{Ablation (Study C)}

\begin{table}[t]
\centering
\caption{\textsc{Simulation-based (synthetic)} Study~C: one component removed at a time, on the same 100 scenarios (policy level). Each row is a simulated configuration, not a run of the implementation; the rows for $\Phi$, memory and recovery correspond to switches that exist in the code (\texttt{verification\_mode}, \texttt{enable\_memory}/\texttt{enable\_rag}, \texttt{enable\_recovery}). $^{\S}$Simulated removal of all stored evidence (screenshots and records); the implemented switch \texttt{enable\_evidence} removes only the per-action execution records. Decisions are unchanged because $\Phi$ reads the live page, not the stored evidence (code behaviour). Evid.: share of episodes with audit evidence stored. MC: means over 1000 regenerated sets.}
\label{tab:sim-ablation}
\scriptsize
\setlength{\tabcolsep}{2pt}
\begin{tabular}{@{}lccccccc@{}}
\toprule
Configuration & TSR & FCR & FRR & F1 & Evid. & MC TSR & MC FCR \\
\midrule
Full & 37/100 [0.28,0.47] & .339 & .213 & .718 & 0.89 & 0.44 & 0.25 \\
-- stored evidence$^{\S}$ & 37/100 [0.28,0.47] & .339 & .213 & .718 & 0.00 & 0.44 & 0.25 \\
-- $\Phi$ & 43/100 [0.34,0.53] & .394 & .065 & .735 & 0.89 & 0.45 & 0.36 \\
-- memory, RAG & 38/100 [0.29,0.48] & .345 & .191 & .724 & 0.89 & 0.44 & 0.26 \\
-- recovery & 30/100 [0.22,0.40] & .348 & .250 & .698 & 0.90 & 0.38 & 0.25 \\
\bottomrule
\end{tabular}
\end{table}


Table~\ref{tab:sim-ablation} removes one component at a time. Removing $\Phi$ reproduces the no-verification row. In the simulation, disabling the recovery engine, with the retries before it still in place, lowers synthetic TSR from 0.44 to 0.41 (3 successes lost, none gained; $p=0.25$) and, across replications, from 0.47 to 0.45; the size of this effect is set by the recovery probabilities of Study~E, which rest on assumed per-attempt success rates. Removing memory changes one outcome in each direction and nothing on average. Stored evidence is not ablated: $\Phi$ reads the live page, so stored evidence serves the audit trail and does not change decisions.

\subsection{Memory and Retrieval (Study D; H3)}

\begin{table}[t]
\centering
\caption{Study~D, memory and retrieval. Upper block ACTUAL (real store, synthetic records); lower block \textsc{Simulation-based (synthetic)}. P(misleading): probability that the retrieved context misleads the planner; for the store it is set to the measured share of irrelevant top-1 records (DERIVED; that this share is the probability of misleading the planner is ASSUMED); for the filter it is an assumed input, so any effect of the filter follows from that assumption. With 44 repeated workflows the study cannot detect an effect of the assumed size (gain 0.08, loss 0.10). $^{\ddagger}$Not implemented.}
\label{tab:sim-memory}
\scriptsize
\setlength{\tabcolsep}{3pt}
\begin{tabular}{@{}lcccc@{}}
\toprule
\multicolumn{5}{@{}l}{\emph{ACTUAL retrieval, real store, 22 synthetic records; 9 queries with a relevant record}} \\
Precision@3 / Recall@3 / MRR & 0.52 & 1.00 & 0.93 &  \\
Top-1 irrelevant or failed (all 10 queries) & 3/10 &  &  &  \\
\midrule
\multicolumn{5}{@{}l}{\emph{SYNTHETIC repeated workflows (n=44)}} \\
Configuration & TSR & FCR & P(misleading) &  \\
Memory disabled & 19/44 [0.30,0.58] & .387 & -- &  \\
Memory enabled & 19/44 [0.30,0.58] & .367 & 0.30 (der.) &  \\
Relevance-filtered$^{\ddagger}$ & 20/44 [0.32,0.60] & .355 & 0.10 (assumed) &  \\
\bottomrule
\end{tabular}
\end{table}


On the synthetic memory corpus the real store returned a relevant record in the top three for each of the 9 queries that had one (recall@3 1.00, MRR 0.93 over those 9; \textsc{actual}). It has no relevance threshold, however: a query with no relevant memory still returned three records, and for 3 of 10 queries the top record was irrelevant or a \emph{failed} earlier attempt (Table~\ref{tab:sim-memory}). Relevance is labelled by task group and is loose (a search for one place counts as relevant to a search for another), so recall is optimistic, and precision@3 is bounded by the number of relevant records per query; the measured 0.52 is the largest value attainable on this corpus. With a misleading rate set to the measured share of irrelevant top-1 records (3/10; that this share equals the probability of misleading the planner is assumed), the simulated repeated workflows succeed 19 times of 44 with and without memory, so H3 is not supported. With 44 workflows and an assumed gain of 0.08 the study could not have detected an effect of that size, so it is uninformative about H3 rather than evidence against it. The proposed relevance filter (20 of 44) changes the outcome only through its assumed misleading rate.

\subsection{Recovery (Study E; H4)}

\begin{table}[t]
\centering
\caption{\textsc{Simulation-based (synthetic)} Study~E: 120 failure scenarios, 20 per mode; 27 drawn as permanent (25\% expected), none of them a missing confirmation. Entry paths and strategy sequences follow the code: action errors are retried three times inline before the engine is called; navigation failures and requests for help go to the engine directly; a missing confirmation raises no error and is handled only by rejection and continuation. Per-attempt success rates are ASSUMED (PROPOSED for the guided variant), so differences between rows follow from them. Detected: failure noticed; FC: undetected missing confirmations reported complete; Unnec.: attempts on permanent failures; Att.: mean attempts. $^{\ddagger}$Not implemented.}
\label{tab:sim-recovery}
\scriptsize
\setlength{\tabcolsep}{2pt}
\begin{tabular}{@{}lccccc@{}}
\toprule
Recovery & Detected & Recovered & FC & Unnec. & Att. \\
\midrule
Inline retries only & 111/120 [0.86,0.96] & 29/93 [0.23,0.41] & 9 & 51 & 1.28 \\
Retry + engine (impl.) & 111/120 [0.86,0.96] & 67/93 [0.62,0.80] & 9 & 162 & 3.12 \\
Evidence-guided$^{\ddagger}$ & 118/120 [0.94,1.00] & 87/93 [0.87,0.97] & 2 & 162 & 3.17 \\
\midrule
\multicolumn{6}{@{}l}{\emph{Recovered / recoverable by mode (inline only; retry + engine; guided)}} \\
\multicolumn{6}{@{}l}{nav. 0; 14; 16 / 16, time-out 9; 12; 13 / 14, missing el. 9; 13; 14 / 15} \\
\multicolumn{6}{@{}l}{invalid sub. 9; 13; 14 / 14, unexp. state 0; 13; 13 / 14, missing conf. 2; 2; 17 / 20} \\
\bottomrule
\end{tabular}
\end{table}


The real recovery engine classifies a navigation error containing \texttt{net::ERR\_} as transient, treats an invalid submission as a verification failure and assigns it re-planning, and does not recognize a missing confirmation, which it classifies as unknown (\textsc{actual}). In the code, action errors, including a planner's request for help, are first retried three times inline and only then passed to the engine; a failed initial navigation is retried three times by the browser executor before it reaches the engine; and a missing confirmation raises no error at all. Simulated along these paths with assumed per-attempt success rates (Table~\ref{tab:sim-recovery}), the retries before the engine recover 51 of 93 synthetic recoverable failures, and adding the engine raises this to 68 of 93. Missing confirmations are recovered only by continuation after a rejection (2 of 20 with or without the engine), and 9 of them end as false completions. The proposed evidence-guided recovery, which checks an explicit confirmation postcondition, recovers 87 and leaves 2 false completions (19 scenarios recovered only by it, none only by the engine); this is a consequence of its assumed success rates, not evidence for H4. In the simulation the engine spends 174 attempts on permanent failures, the guided variant 162.

\subsection{Verification Overhead (Study F; H5)}

\begin{table}[t]
\centering
\caption{Study~F, verification overhead. Upper block ACTUAL; lower block \textsc{Simulation-based (synthetic)}, built from the ACTUAL component samples and an ASSUMED inference time; the two blocks are not combined into one claim. Inference dominates every synthetic episode (84\% of time under $\Phi$); evidence screenshots take 2.5\%. $^{\ddagger}$Not implemented.}
\label{tab:sim-overhead}
\scriptsize
\setlength{\tabcolsep}{3pt}
\begin{tabular}{@{}lccc@{}}
\toprule
\multicolumn{4}{@{}l}{\emph{ACTUAL component latency (ms), local sites, 30 repetitions}} \\
Operation & median & p95 &  \\
navigate & 17.1 & 24.9 &  \\
click + load & 65.2 & 106.2 &  \\
fill & 12.5 & 17.4 &  \\
DOM extraction & 76.4 & 105.6 &  \\
screenshot (evidence) & 61.3 & 93.6 &  \\
$\Phi$ on a site page & 5.2 & 6.9 &  \\
$\Phi$ on gate states (83) & 26.8 & 30.3 &  \\
\midrule
\multicolumn{4}{@{}l}{\emph{SYNTHETIC episode time (s), inference median 3.0 s assumed}} \\
Configuration & median & p95 & verif.\ share \\
No verification & 17.0 & 34.1 & 0.0\% \\
Evidence-gated ($\Phi$) & 17.0 & 34.2 & 0.2\% \\
LLM self-verification & 21.7 & 38.7 & 19.6\% \\
$\Phi$ + model check$^{\ddagger}$ & 21.7 & 38.7 & 19.8\% \\
\bottomrule
\end{tabular}
\end{table}


$\Phi$ takes a median of 5.2\,ms on an ordinary page of the self-hosted sites and 26.8\,ms on the constructed gate states; a screenshot 61\,ms and DOM extraction 76\,ms (\textsc{actual}, Table~\ref{tab:sim-overhead}). In the simulation, with an assumed median of 3\,s per model call on a CPU, inference takes 80\% of an episode and the checks of $\Phi$ 0.1\%. The mean episode is nevertheless 2.6\% longer with $\Phi$ (20.2 against 19.7\,s), because each rejected claim costs one further planning step (planner call, DOM extraction, click, three screenshots; assumed), and $\Phi$ times are resampled from the slower gate-state measurements; a model self-check adds 14\% (22.4\,s; a cost-only figure, since the simulated self-check accepts every claim). At 1.5\,s or 6\,s per call the median episode under $\Phi$ takes 10.0\,s or 36.8\,s. H5 is consistent with the simulation: the rule itself is cheap, and its cost is in false rejections and the re-planning that rejections cause.

\subsection{Robustness and Sensitivity (Study H)}

\begin{table}[t]
\centering
\caption{\textsc{Simulation-based (synthetic)} Study~H: robustness of the evidence-gated configuration (means over 300 regenerated sets per row, drawn from a random stream separate from the 1{,}000-set Monte Carlo of Table~\ref{tab:sim-tests}, so the unperturbed rows differ slightly from it). Without verification FCR stays at 0.36 at every continuation rate.}
\label{tab:sim-sens}
\scriptsize
\setlength{\tabcolsep}{3pt}
\begin{tabular}{@{}lccc@{}}
\toprule
Varied factor ($\Phi$) & FCR & FRR & TSR \\
\midrule
$\Phi$ detection $\times0.5$ & 0.30 & 0.14 & 0.46 \\
$\Phi$ detection $\times1.0$ & 0.25 & 0.14 & 0.47 \\
$\Phi$ detection $\times1.5$ & 0.19 & 0.13 & 0.49 \\
missing evidence 0\% & 0.23 & 0.09 & 0.52 \\
missing evidence 8\% & 0.25 & 0.14 & 0.47 \\
missing evidence 20\% & 0.29 & 0.22 & 0.39 \\
action failure rate 0.1 & 0.25 & 0.14 & 0.50 \\
action failure rate 0.25 & 0.25 & 0.14 & 0.47 \\
action failure rate 0.4 & 0.25 & 0.14 & 0.44 \\
finish after rejection 0.1 & 0.26 & 0.14 & 0.45 \\
finish after rejection 0.35 & 0.25 & 0.14 & 0.47 \\
finish after rejection 0.6 & 0.24 & 0.13 & 0.50 \\
\bottomrule
\end{tabular}
\end{table}


The advantage of verification depends on how often the planner claims completion prematurely, which is unmeasured (Fig.~\ref{fig:sim-sens}): at $\pi=0.3$ the FCR of $\Phi$ is 0.14 against 0.21 without verification, at $\pi=0.95$ 0.32 against 0.45, while FRR stays near 0.14 against 0.08. Halving the derived detection rates of $\Phi$ raises its FCR from 0.25 to 0.30, about halfway to no verification (0.36); missing evidence raises FRR (0.09 to 0.22) because the simulated $\Phi$ rejects a claim on an unreadable page with assumed probability 0.8; more action failures lower TSR; and a higher rate of finishing after a rejection raises the TSR of $\Phi$ from 0.45 to 0.50 (Table~\ref{tab:sim-sens}).
